\subsection{Organization and novelty boundary}

Sections~\ref{sec:relations} and \ref{sec:fitting} establish the
relation model and Fitting presentation.  The intrinsic filtration
leads directly to Section~\ref{sec:intrinsic-web-reconstruction}:
the deepest cone, the universal coefficient readout, its functoriality,
and the exterior-support theorem form the proof of reconstruction.
The polarized K3, moduli, and classical inverse-data comparisons then
locate the result among the associated geometric constructions.

Appendix~\ref{app:boundary-overview} states the full boundary and
family theorems.  The subsequent appendices retain every coefficient
table, collision calculation, determinant-completion proof, relative
primary-filtration argument and packet-descent construction.  The
separating pencils, weighted models and exact certificates are also
included.  These results are not hypotheses of the support argument;
in particular it does not require a full embedded-primary atlas of
\(W_3,W_4\) at the highest coranks.

Deformation to the normal cone, flattening stratification, Schur
functors, and standard monomial straightening are classical tools.
The new assertions are the specific residual Fitting algebra, the
universal determinant-completion mechanism applied to it, effective
Pieri nonvanishing for the corank-four membership, the generic and
primary corank-two boundary tables, their specialization laws, and
the reconstruction of the polarized K3 and of every smooth Jacobian
web from the intrinsic nonreduced scheme.  Ballico's 1993 paper remains cited as a documentary boundary;
no theorem-level nonanticipation claim is made without its complete
text.

The relative theorem is positioned in
Section~\ref{sec:relative-tool-positioning}: its role is simultaneous
control of a finite exact diagram, rather than a new generic-freeness
principle.  Section~\ref{sec:primary-table-scope} separates complete
corank-two and benchmark primary decompositions from the global
support/depth laws.  This does not alter the full-scheme inverse
statement, which reads the maximal-minor space without requiring an
embedded-primary decomposition of every higher residual layer.
